% ============================================================================
%  Actividad 4 · Guía de estilos del producto digital — once secciones
%
%  Propiedad de US-UX-09. Las dos secciones puntuadas restantes (conclusiones
%  y referencias) viven en a4_06_cierre.tex, de la misma historia.
%
%  Regla de este archivo: ningún hexadecimal se teclea como instrucción de
%  color. Todo color del informe entra por \laminaPaleta, \laminaContraste,
%  \laminaTipografia y \laminaReticula, que emite
%  docs/entregables/generar_tokens_a4.py desde los 37 \definecolor de
%  estilo/uxdoc.sty, o por nombre de color de esa misma hoja.
%
%  Segunda regla, heredada del bloque del tema institucional que este archivo
%  absorbió: ningún valor de color se escribe con almohadilla.
%  scripts/verificar_tokens_a4.sh falla si encuentra un hexadecimal con
%  almohadilla en cualquier contenido/a4_*.tex. Los valores del sistema de
%  tokens del portal se publican como texto monoespaciado, que es lo que son
%  aquí: el dato de procedencia de un token, no una instrucción de color.
%
%  Versión de la guía: v1.0. La cadena impresa en el PDF NO se teclea aquí,
%  sale de \versionguia y \fechaguia, que emite el generador. Esta línea de
%  comentario existe porque scripts/verificar_tokens_a4.sh exige que el
%  documento nombre la misma versión que el bloque @theme, la paleta tipada y
%  el manifiesto; el verificador falla si las cuatro dejan de coincidir.
% ============================================================================

\begin{uxdestacado}
{\uxhead\fontsize{13.5}{16}\selectfont\color{uxnavy}Guía de estilos de Karisma Data · \versionguia\ · \fechaguia}\par
\vspace{6pt}
{\color{uxamber}\rule{60pt}{2.2pt}}\par
\vspace{8pt}
La arquitectura que la Actividad 3 validó ---cuatro categorías, dieciséis ramas y nueve facetas transversales--- llegó a esta actividad como contrato de navegación, y esta guía es la otra mitad del contrato: la que fija cómo se ve y cómo se comporta cada pieza que ocupa esas rutas. Las once secciones siguientes cubren identidad ---la del informe y la del portal---, tipografía, el sistema de tokens con sus dos temas y sus dos modos, componentes, iconografía, imágenes, retícula, movimiento, accesibilidad, lenguaje y versionado.

\vspace{6pt}
La guía se deriva de un solo archivo y en un solo sentido. Los 37 nombres de color viven en \texttt{estilo/uxdoc.sty}; el generador \texttt{generar\_tokens\_a4.py} los lee y emite cuatro salidas: el bloque \texttt{@theme} que la aplicación consume, la paleta tipada que la ruta \texttt{/guia} imprime, las láminas de este documento y un manifiesto legible por máquina. La consecuencia práctica es que la paleta que este PDF imprime y la que el navegador pinta no divergen: son el mismo dato leído dos veces. Una guía de estilos que no está enlazada al código envejece el día que se entrega, y esa es la razón de la cadena, no una comodidad de implementación.

\vspace{6pt}
La segunda regla de alcance es igual de estricta: \textbf{esta guía no documenta ningún componente que el producto no renderice}. Las \textbf{once} láminas vivas de la ruta \texttt{/guia} ---paleta, tipografía, botones, campos, tablas, tarjetas, chasis, iconos, linaje, accesibilidad y tarjetas de llamada a herramienta--- son la fuente de las secciones 3, 4 y 5. La lámina de chasis pertenece al eje de tema y publica los cuatro tokens de la barra lateral, los tres estados de certificación y la superficie contenida en sus cinco canales. La anatomía sigue la del material de apoyo del curso (Ramírez Mejía, 2023) y la amplía en las seis secciones que ese material no cubre.
\end{uxdestacado}

% ============================================================================
\section{Identidad de la marca}

\textbf{Karisma Data · Portal Centralizado de Datos Financieros.} La marca nombra lo que el producto hace: reunir fuentes dispersas en un único dato con autoridad. El símbolo representa esa convergencia y el logotipo compuesto se arma con el símbolo más el nombre en Lexend Deca, de modo que el texto de la marca es siempre vectorial y nunca una imagen rasterizada dentro de otra imagen. La sección documenta los \textbf{dos sistemas de marca} que el proyecto mantiene: el del informe impreso, con su construcción de portada, y el del portal, con su geometría medida sobre la página normativa del archivo de diseño.

\subsection{Construcción del logotipo compuesto}

La construcción está fijada en la macro \texttt{\textbackslash uxportada} de \texttt{estilo/uxdoc.sty} y se aplica idéntica en las cuatro entregas del curso: símbolo a 2{,}0 cm de ancho, nombre debajo en Lexend Deca a 23 pt sobre el azul institucional, y descriptor \textit{Portal Centralizado de Datos Financieros} en cursiva sobre el neutro secundario. La portada de este documento es la aplicación de esa regla, no una composición aparte.

\subsection{Área de resguardo y tamaño mínimo}

\begin{uxtabla}{|L{4.0cm}|L{3.0cm}|Y|}{Reglas de resguardo y de tamaño mínimo de la marca}
  \uxheadrow \thd{Regla} & \thd{Medida} & \thd{Motivo} \\
  Área de resguardo & Una altura de símbolo por cada lado & Ningún elemento entra en esa banda: ni filete, ni texto, ni borde de tarjeta \\
  Tamaño mínimo del símbolo solo & 16 px de alto & Por debajo, los trazos de convergencia se funden y el símbolo se lee como una mancha \\
  Tamaño mínimo del compuesto & 96 px de ancho & Es el ancho al que el nombre sigue siendo legible junto al símbolo \\
  Proporción & Bloqueada & El compuesto escala en un solo gesto; alto y ancho nunca se ajustan por separado \\
\end{uxtabla}

\subsection{Versiones permitidas}

Tres, y ninguna más. \textbf{Positiva} sobre la superficie clara del sistema, que es el uso por omisión. \textbf{Negativa} sobre la primaria 700 y sobre la primaria 900, que son los dos fondos oscuros del producto: la barra lateral usa exactamente esa combinación y la matriz de la sección 3 la respalda con 8{,}39:1 y 13{,}66:1. \textbf{Símbolo solo}, reservado al favicon y al avatar, donde no cabe el nombre.

\subsection{Usos incorrectos}

\begin{uxtabla}{|L{5.2cm}|Y|}{Los cuatro usos incorrectos de la marca}
  \uxheadrow \thd{Uso incorrecto} & \thd{Qué rompe} \\
  Deformar la proporción & Cambia el peso relativo del símbolo frente al nombre y delata una composición manual \\
  Recolorear fuera de la paleta & Rompe la única fuente de color: cualquier tono que no salga de \texttt{uxdoc.sty} es un color escrito a mano \\
  Colocar sobre fotografía sin capa de contraste & El fondo variable impide garantizar la razón mínima, que es justo lo que la sección 3 verifica por cálculo \\
  Añadir sombra, contorno o degradado & El sistema declara tres niveles de sombra y ninguno se aplica a la marca \\
\end{uxtabla}

\subsection{Atributos de marca}

Cuatro, en orden de prioridad: \textbf{confiable · preciso · sobrio · rápido}. El orden es operativo y no decorativo: cuando dos opciones de diseño empatan, gana la que refuerza el primero. Es la regla que decidió que la cifra de una respuesta del asistente aparezca solo después de la tarjeta que la consultó, y que el ámbar de acento quedara reducido a un papel decorativo en cuanto el cálculo mostró que no alcanzaba el contraste de un elemento informativo.

\subsection{La marca del portal, medida sobre la página normativa}

La marca \textbf{del portal} se rige por la página normativa del archivo de diseño, titulada
\enquote{Uso del logotipo}: un símbolo en forma de \textbf{K sobre una teja de esquinas
redondeadas}, con tres variantes y sus reglas de aplicación. Se dibuja como gráfico vectorial en el
propio portal y no se genera como imagen, por las reglas de esa misma página: prohíben deformar y
recolorear, y un archivo de mapa de bits producido aparte sería una reinterpretación con vida
propia. Al dibujarla, la geometría no se estimó a ojo: se midió sobre la lámina exportada a 600
puntos por pulgada y se expresó en proporción del lado de la teja, que es lo que permite escalarla
sin volver a decidir nada.

La teja mide \textbf{112 píxeles de diseño} y es cuadrada. Todo lo demás va en porcentaje de ese
lado.

\begin{uxtabla}{|L{4.4cm}|L{3.6cm}|Y|}{Geometría del símbolo, medida sobre la página normativa del archivo}
  \uxheadrow \thd{Elemento} & \thd{Proporción del lado} & \thd{Nota} \\
  Radio de la teja & 6{,}01 \% & Esquina redondeada del contenedor \\
  Radio de las barras & 2{,}47 \% & La barra ámbar usa un radio ligeramente menor \\
  Ancho de columna & 14{,}06 \% & Las tres columnas son iguales \\
  Canal entre columnas & 7{,}40 \% & Derivado, e idéntico en los dos huecos \\
  Posición de las columnas & 17{,}81 \%, 39{,}27 \% y 60{,}73 \% & Asta, brazo alto y brazo corto \\
  Borde superior & 19{,}64 \% & Único para las tres \\
  Alturas & 60{,}73 \%, 39{,}27 \% y 23{,}18 \% & Descendentes, de izquierda a derecha \\
  Barra de base & 35{,}52 \% por 14{,}16 \% & Situada en 39{,}27 \% y 66{,}21 \% \\
\end{uxtabla}

Dos cifras de la barra de base \textbf{no} son la medición cruda. Su borde derecho y su línea de
base salían dos píxeles ---a 600 puntos por pulgada--- de los de la tercera columna y del asta, y
eso es el suavizado del trazado y no una decisión de diseño, así que se igualaron. El margen
vertical queda simétrico: 19{,}64 \% arriba y 19{,}63 \% abajo.

\subsection{Las tres reglas de la página normativa y las tres variantes de la marca}

\begin{uxtabla}{|L{3.6cm}|L{3.8cm}|Y|}{Las reglas de aplicación de la marca y cómo se implementan}
  \uxheadrow \thd{Regla del archivo} & \thd{Enunciado} & \thd{Cómo la cumple el portal} \\
  Área de protección & Espacio libre mínimo alrededor del símbolo equivalente a media altura de la K & Reserva de 16 píxeles junto al símbolo de 32, que es exactamente esa mitad \\
  Tamaño mínimo & Símbolo digital de 32 píxeles; marca completa de 120 píxeles de ancho & El símbolo se dibuja a 32 píxeles, que es a la vez su tamaño y su piso; el nombre solo aparece cuando la marca completa alcanza sus 120 píxeles de ancho, es decir a partir de 768 píxeles de pantalla \\
  Usos incorrectos & No deformar, rotar, recolorear, separar elementos ni añadir efectos & La proporción viaja en el propio lienzo del gráfico vectorial y los cuatro colores de la marca están fijados con su valor, no con un token del tema \\
\end{uxtabla}

La última implementación merece explicación, porque es la única excepción a la regla de que todo
color del portal sale de un token: \textbf{la marca no se repinta con el tema}. El color de acción
es verde azulado bajo el tema institucional y tinta bajo el de omisión, así que pintar la teja con
ese token habría cambiado la marca al cambiar de tema, que es exactamente lo que la página normativa
prohíbe. Una prueba automática vigila la excepción al revés: falla si el dibujo de la marca llega a
contener una referencia a una variable de color del sistema.

\begin{uxtabla}{|L{3.4cm}|L{3.4cm}|L{3.4cm}|Y|}{Las tres variantes de la marca y sus colores medidos}
  \uxheadrow \thd{Variante} & \thd{Teja} & \thd{Barras claras} & \thd{Barra de base} \\
  Principal & \texttt{086B70} (Acción) & \texttt{FFFFFF} (Superficie) & \texttt{B97812} (Atención) \\
  Inverso & \texttt{086B70} & \texttt{FFFFFF} & \texttt{B97812} \\
  Monocromático & \texttt{102A43} (Navegación) & \texttt{FFFFFF} & \texttt{FFFFFF} \\
\end{uxtabla}

Principal e inverso son \textbf{el mismo dibujo}: se midieron las tres tejas por separado y sus
histogramas coinciden. Lo que cambia en la lámina es el fondo de la tarjeta sobre la que se
presentan, no el relleno del símbolo. Para que la variante no fuera decorativa, gobierna además el
color del nombre en la marca completa: hereda la tinta de quien la monta en la principal, blanco en
la inversa y azul profundo en la monocromática.

Nótese que el ámbar de la marca es el valor literal del archivo, \texttt{B97812}, y no el token
\texttt{aviso}, que es ese mismo ámbar oscurecido un 12 \%. La diferencia no es una incoherencia:
un color de identidad y un color que informa se rigen por reglas distintas, y el listón de 4{,}5:1
aplica al segundo. El ámbar de la marca ocupa aquí un área gráfica de tres píxeles de alto dentro de
un símbolo, no una etiqueta que alguien tenga que leer.

\subsection{La marca del informe y la marca del portal}

Las cinco subsecciones de apertura gobiernan la marca \textbf{de este informe}: la construcción de
la portada, su área de resguardo y sus tamaños mínimos, que son medidas de tinta sobre papel. Las
dos anteriores gobiernan la marca \textbf{del portal}, cuyas reglas salen de la página normativa
del archivo de diseño. Los dos sistemas están separados a propósito y sus cifras no coinciden ---32
y 120 píxeles en pantalla, 16 y 96 sobre papel---, porque un mínimo legible en pantalla y uno
legible impreso no son el mismo número. Del portal al informe viaja contenido, nunca formato.

\subsection{Convivencia con la marca institucional}

Las dos marcas conviven en la portada de los entregables del curso ---el escudo del Tecnológico de Monterrey arriba y la marca del producto en el bloque central--- y no conviven en la interfaz. Dentro de la aplicación aparece únicamente la marca del producto, que es \enquote{Karisma Data} en los dos idiomas. La razón es de propiedad: el portal es un caso de estudio académico y la marca institucional acredita el trabajo, no el producto.

% ============================================================================
\section{Tipografía}

Dos familias, servidas desde el propio origen por el módulo de fuentes de la aplicación y cargadas con \texttt{fontspec} en este documento. La misma pareja compone la interfaz y el PDF, así que el espécimen impreso más abajo es el tipo real del producto y no una aproximación.

\subsection{Las dos familias y su justificación}

\textbf{Titulares: Lexend Deca.} Elegida por su legibilidad en tamaños grandes y por un carácter corporativo sin rigidez. Se distribuye únicamente en peso Regular.

\textbf{Cuerpo: Fira Sans.} Humanista, con altura de x generosa y buenas cifras. Rinde bien entre 12 y 14 px, que es la escala real de un producto denso de escritorio.

\subsection{Espécimen y escala de nueve roles}

\laminaTipografia

\subsection{El peso 400 es una decisión del sistema}

Los nueve roles llevan peso 400 y la jerarquía se construye con tamaño, color y espacio en blanco. No es una carencia: Lexend Deca solo se distribuye en Regular y la configuración de fuentes descarga esa única variante, de modo que un peso 600 en el bloque \texttt{@theme} produciría negrita simulada en la pantalla y negrita real en el PDF. Las capturas de la guía dejarían de coincidir con el documento, que es el defecto exacto que esta cadena existe para impedir. El rol \texttt{text-etiqueta} se distingue con espaciado entre letras en lugar de con peso.

\subsection{Cifras tabulares, obligatorias}

Toda columna numérica, todo importe y todo contador se componen con cifras tabulares, el valor \texttt{tabular-nums} de la propiedad que fija el ancho de los dígitos. Sin eso las columnas bailan al actualizarse y un portal financiero pierde credibilidad en el primer parpadeo. La regla alcanza al rol \texttt{text-dato}, que es el de la cifra de tarjeta KPI y el de la celda numérica.

\subsection{La tensión 14/16 px, resuelta y no ignorada}

La recomendación general de 16 px de cuerpo está pensada para lectura móvil; la ficha de estilo que corresponde a un producto de tablero denso fija entre 12 y 14 px. El sistema resuelve la tensión declarando las dos reglas: \texttt{text-cuerpo} a 14 px es el texto por omisión en escritorio, y los campos de formulario por debajo del quiebre de 768 px suben a 16 px para evitar el acercamiento automático de los navegadores móviles. El rol \texttt{text-cuerpo-amplio}, a 16 px con interlínea de 26 px, cubre la prosa larga: ayuda, respuesta del asistente y aviso de alcance.

\subsection{Medida de línea}

Entre 60 y 75 caracteres en todo bloque de prosa. Las tablas y los paneles quedan exentos por naturaleza: su medida la fija la retícula de la sección 7.

% ============================================================================
\section{Especificación del sistema de tokens y paleta institucional}

Cuatro familias de cinco tonos, siete neutros, cuatro semánticos y seis colores de serie categórica: \textbf{37 nombres de color}, todos declarados en \texttt{estilo/uxdoc.sty}. Once de ellos son las anclas de la Actividad 1 y conservan su valor byte a byte, porque los documentos de A1, A2 y A3 ya compilan contra ellas; los veintiséis restantes se añadieron en un bloque propio, veinte con \texttt{\textbackslash definecolor} y seis con \texttt{\textbackslash colorlet}. Esos seis son reúsos explícitos ---el aviso apunta al acento 700, el informativo a la primaria 500, tres series a colores de marca--- y apuntar en vez de copiar es lo que impide que un token y su origen se separen con el tiempo.

El capítulo especifica además los \textbf{dos temas} y los \textbf{dos modos} que el sistema
entrega: la paleta de omisión, con su matriz de cuarenta y cuatro pares, y la institucional, con la
procedencia de cada uno de sus veintiocho tokens, su matriz de contraste en los dos modos, su
separación bajo dicromacia y su medición sobre el portal en ejecución. Las dos se emiten con el
mismo generador y se juzgan contra el mismo umbral, que es lo que permite publicarlas en una sola
especificación en lugar de en dos.

\subsection{Paleta completa}

\laminaPaleta

\subsection{La regla que gobierna la paleta}

\textbf{Ningún suelo es negro puro, y una sola superficie es blanco puro.} El suelo del sistema es \texttt{ground} y la tinta más oscura es \texttt{corriente-pleno}. En el tema de omisión valen \texttt{F4F6F9} y \texttt{14171D}; en el institucional, \texttt{FFFFFF} y \texttt{102A43}. Los cuatro suelos y las cuatro tintas se mantienen dentro del canal ---ninguno toca el negro--- con \textbf{una excepción que el sistema declara en lugar de esconder}: el suelo claro del tema institucional \emph{sí} es blanco puro. No es un descuido ni una concesión: \emph{Superficie} del archivo de identidad es blanco puro, y este tema existe para llevar ese archivo, no para corregirlo. El tema de omisión conserva su \texttt{F4F6F9} y sigue siendo la superficie que la regla describe, de modo que el producto nunca queda sin un suelo conforme. La consecuencia es medible y la sección 3.4 la documenta: la matriz calcula el contraste contra el suelo \textbf{real de cada tema} y no contra un blanco supuesto, que es la razón por la que las razones de un tema no se pueden leer como las del otro.

\subsection{Matriz de contraste verificada por cálculo}

Las razones de contraste no se estimaron a ojo ni se copiaron de una tabla ajena: el generador las calcula para \textbf{44 pares} con la fórmula de WCAG 2.x ---luminancia relativa de cada color y cociente $(L_1+0{,}05)/(L_2+0{,}05)$--- y emite el veredicto contra los umbrales de la norma: 4{,}5:1 para texto normal, 3:1 para texto grande y para elementos gráficos, 7:1 para el nivel AAA (W3C, 2023). La lámina siguiente es esa matriz completa, con el veredicto de cada par, la tabla de los cuatro defectos que encontró y las cinco reglas que salieron de ella.

\laminaContraste

\subsection{Los cuatro defectos que la matriz corrigió}

El cálculo no confirmó el sistema: lo corrigió en cuatro puntos, y los cuatro estaban escritos como reglas buenas en la planeación de la semana. Conviene decirlo con los números a la vista, porque es la evidencia de que la verificación por cálculo hizo trabajo real y no ceremonia.

\begin{enumerate}[leftmargin=*,itemsep=5pt,topsep=6pt]
\item \textbf{El número del ámbar estaba mal.} La planeación citaba 2{,}6:1 para el texto claro sobre el acento 500. El cálculo da \textbf{2{,}80:1} con blanco puro y \textbf{2{,}68:1} con la superficie que el producto usa de verdad. El veredicto no cambió ---los dos fallan---, pero un número equivocado impreso con el rótulo \enquote{verificado} es peor que ningún número. La regla sustituta es que, si el acento 500 lleva texto, ese texto es \texttt{ink}: 5{,}22:1, nivel AA.
\item \textbf{El texto de aviso fallaba AA.} El acento 700 sobre el acento 100 se daba por bueno y da \textbf{3{,}96:1}, por debajo del 4{,}5:1 de texto normal. El texto de aviso pasó al \textbf{acento 900} sobre el mismo fondo: \textbf{7{,}58:1}, nivel AAA, sin cambiar el fondo ni el tamaño de letra.
\item \textbf{El ámbar 500 no llega a elemento gráfico.} Sobre la superficie da \textbf{2{,}68:1} y el umbral de un límite de componente es 3:1. El acento 500 quedó \textbf{solo decorativo} ---filete y marca--- y todo ámbar que informa usa el acento 700, con \textbf{4{,}40:1}.
\item \textbf{El borde de campo era invisible.} El filete claro \texttt{line} sobre la superficie da \textbf{1{,}42:1}: el campo de formulario se veía como un rectángulo sin borde. El borde de campo pasó al neutro \texttt{line-strong}, con \textbf{4{,}55:1}, y el filete claro quedó para separadores decorativos.
\end{enumerate}

Los cuatro se corrigieron \textbf{en los tokens y no en la prosa}, que es la diferencia entre documentar un defecto y arreglarlo. El efecto colateral se midió en vez de suponerse: el token \texttt{accent-text} cambió de acento 700 a acento 900 y su único consumidor de entonces ---la insignia del botón de prototipo sobre la superficie alterna--- pasó de 4{,}12:1, que falla AA, a \textbf{7{,}89:1}, nivel AAA, sin tocar una línea del componente.

\subsection{Las cinco reglas derivadas}

Las cinco reglas que cierran la lámina anterior no son recomendaciones: son la forma normativa de los cuatro hallazgos, y cada una tiene un token que la hace cumplir. Quien escribe una pantalla no elige entre \texttt{line} y \texttt{line-strong} por criterio propio; usa el token de borde de campo, que ya apunta al valor que pasa el umbral. Ese es el mecanismo que evita que una regla escrita se pierda entre la guía y el código.

\subsection{Paleta categórica de series}

Una gráfica con seis series en tonos del mismo azul es ilegible, así que la paleta de series es distinta de la de marca aunque reúse tres de sus colores. Las seis quedan por encima de 3:1 sobre la superficie ---la peor es 4{,}40:1--- y ninguna depende del color para distinguirse: cada serie lleva además marcador y patrón de línea propios.

\begin{uxtabla}{|L{2.1cm}|L{3.0cm}|L{2.6cm}|Y|}{Paleta categórica de series: color, marcador y patrón de línea}
  \uxheadrow \thd{Serie} & \thd{Marcador} & \thd{Línea} & \thd{Razón sobre la superficie} \\
  \texttt{serie-1} & Círculo & Sólida & 4{,}94:1 \\
  \texttt{serie-2} & Triángulo hacia arriba & Discontinua & 4{,}40:1 \\
  \texttt{serie-3} & Cuadrado & Punteada & 6{,}81:1 \\
  \texttt{serie-4} & Rombo & Guion y punto & 5{,}45:1 \\
  \texttt{serie-5} & Triángulo hacia abajo & Sólida gruesa & 5{,}12:1 \\
  \texttt{serie-6} & Cruz & Discontinua fina & 7{,}54:1 \\
\end{uxtabla}

\subsection{Los semánticos nunca son categorías}

\texttt{danger}, \texttt{warning}, \texttt{info} y \texttt{success} comunican estado y no pertenencia a una categoría. Una serie de gráfica jamás se pinta con ellos, porque el lector leería un juicio donde solo hay una etiqueta. La regla inversa también aplica: un estado nunca se pinta con un color de serie.

\subsection{Los dos temas y los dos modos}

El sistema de diseño del portal declara \textbf{dos temas} y \textbf{dos modos} como capacidades
nativas de su arquitectura, no como una capa añadida sobre una paleta única. El eje de tema declara
exactamente dos cosas ---color y familia tipográfica--- y el eje de modo declara el suelo sobre el
que ese color se lee. Las \textbf{cuatro combinaciones} que resultan están verificadas una por una
con el generador que emite los tokens del portal y medidas después en el navegador. Ninguno de
los cuatro suelos es negro puro, por la regla de la subsección 3.2: casi negro en el tema de
omisión y azul profundo en el institucional.

\textbf{El tema institucional es el tema central del producto.} Su color procede del archivo de
diseño de esta actividad, lleva la identidad de Karisma Data y es el que compone las capturas de la
sección de prototipos. Se elige desde la cabecera del portal, viaja en una cookie y el servidor lo
aplica antes del primer trazado, de modo que el cambio no produce destello.

\textbf{El tema de omisión es una alternativa deliberada y no un residuo.} Es una rampa de alto
contraste sobre neutros, con la cuadrícula modular pintada, y existe porque hay condiciones de uso
en las que un tema de identidad no es la mejor lectura del mismo sistema:

\begin{uxlist}
  \lead{Lectura prolongada de datos densos}{Una jornada de trabajo sobre tablas de treinta y cuatro
    píxeles de fila se sostiene mejor con una rampa neutra: el color queda entonces reservado al
    canal semántico y a las series, y no compite con ellos desde el chasis.}
  \lead{Pantallas de baja calidad y luz ambiente alta}{Un equipo con gama reducida o un puesto junto
    a una ventana aplanan las diferencias de matiz antes que las de luminancia. La rampa de alto
    contraste separa por luminancia, que es la dimensión que sobrevive a las dos condiciones.}
  \lead{Sensibilidad al contraste cromático}{Para quien encuentra fatigante una superficie
    saturada, un tema que separa por luminancia y no por matiz reduce la carga sin retirar
    información: los dos temas cumplen los mismos listones y publican los mismos estados.}
\end{uxlist}

El producto ofrece por tanto \textbf{dos lecturas del mismo sistema} y ninguna de las dos es un
accidente. Lo que ninguna de las dos puede hacer es cambiar la estructura: el eje de tema mueve
color y familia tipográfica, y nada más. El tema de omisión conserva además sus valores sin una sola
diferencia, y esa inmovilidad está fijada por una prueba que compara cada uno de ellos contra su
literal.

\subsection{El octeto del archivo de diseño, con su destino}

El archivo de diseño declara ocho colores con su función escrita al lado. Los ocho ocupan una
ranura del contrato de tokens; ninguno se queda fuera.

\begin{uxtabla}{|L{2.4cm}|C{1.8cm}|L{3.6cm}|Y|}{Los ocho colores que declara el archivo de diseño y su destino en el contrato de tokens}
  \uxheadrow \thd{Nombre en el archivo} & \thd{Valor} & \thd{Función declarada} & \thd{Destino en el contrato de tokens} \\
  Navegación & \texttt{102A43} & Estructura de navegación & \texttt{corriente-pleno} en claro, \texttt{ground-alt} en oscuro y suelo de la barra lateral en los dos modos \\
  Secundario & \texttt{1D4C6E} & Segundo nivel de estructura & \texttt{corriente-medio} en claro y base del canal informativo \\
  Acción & \texttt{086B70} & Acciones y selección & \texttt{accion} y bloque del módulo en curso de la barra lateral \\
  Apoyo & \texttt{15989A} & Acento de apoyo & \texttt{accion-apoyo}: realce, subrayado en curso y filete de foco \\
  Atención & \texttt{B97812} & Atención & \texttt{aviso}, con la divergencia de luminancia declarada abajo \\
  Éxito & \texttt{287A58} & Confirmación & \texttt{ok}, con su valor literal \\
  Error & \texttt{B8443F} & Error & \texttt{error}, con la divergencia de luminancia declarada abajo \\
  Superficie & \texttt{FFFFFF} & Superficie de trabajo & \texttt{ground} del modo claro, con su valor literal \\
\end{uxtabla}

Cinco de los ocho entran con su valor literal, uno ---Navegación--- ocupa tres ranuras distintas
según el modo y el papel, y dos entran corregidos en luminancia por una razón medida. Ese reparto es
la diferencia entre adoptar una paleta y verificarla: el tema institucional toma del archivo todo lo
que el contrato puede alojar, deriva el resto con una regla escrita y mide el resultado antes de
publicarlo.

Los códigos de color de este capítulo se imprimen sin el signo de número que suele precederlos, y no
es un descuido de composición. Una comprobación del repositorio recorre los archivos de contenido de
esta actividad buscando exactamente ese patrón, para impedir que un color se teclee a mano en el
documento en lugar de salir de su generador. Los valores citados aquí describen la procedencia de
una paleta que vive en el sistema de tokens del portal ---otra cadena, con otro generador--- y por
eso aparecen en esta forma, que la comprobación distingue de un color escrito a mano.

\subsection{El chasis, y qué cambia con el tema}

El portal monta \textbf{un solo chasis} para los dos temas: barra lateral con etiquetas de texto,
cabecera fina y franja de alcance, en las diez rutas. Dos chasis habrían sido dos productos y
habrían duplicado cada estado no feliz. Lo que cambia con el tema es el color de esas piezas, y
cambia porque son tokens y no condiciones escritas dentro de un componente.

\begin{uxtabla}{|L{3.2cm}|L{5.0cm}|Y|}{Lo que el tema cambia en el chasis, y lo que deja intacto}
  \uxheadrow \thd{Pieza} & \thd{Tema de omisión} & \thd{Tema institucional} \\
  Fondo del chasis & Retícula modular visible & Sin retícula: el token se pinta de su propio suelo \\
  Barra lateral & Suelo alterno; el módulo en curso se distingue por luminancia y peso & Azul profundo \texttt{102A43}; el módulo en curso es un bloque relleno de acción \\
  Acción primaria & Corriente plena, \texttt{14171D} & Verde azulado \texttt{086B70} \\
  Selección & Tinte neutro & Tinte de acción, \texttt{E6F2F1} \\
  Tarjeta contenida & Filete de un pelo, sin radio & Filete más radio y barra de canal a la izquierda \\
  Familia tipográfica & Lexend Deca en titulares y Fira Sans en cuerpo & Inter en los dos roles \\
\end{uxtabla}

\textbf{Lo que no cambia con el tema}: la estructura, el orden de tabulación, los estados no
felices, la franja de alcance, la densidad de las tablas y las seis series de datos. Un tema que
cambiara la estructura sería otro producto, y el eje del sistema declara exactamente dos cosas
---color y familia tipográfica--- para que nadie pueda meter una tercera sin decirlo.

\subsection{Los veintiocho tokens, con la procedencia de cada valor}

El contrato del sistema declara veintiocho tokens de color y el tema institucional los resuelve todos
en los dos modos. Se reparten en cuatro bloques: los diecisiete de la rampa base y el canal
semántico, los tres de acción, apoyo y selección, los cuatro del chasis y los tres de certificación.
La columna de procedencia distingue tres orígenes: el valor literal del archivo de diseño, una
derivación declarada a partir de uno de sus colores, o el valor heredado del tema de omisión.

\begin{uxtablalarga}{|L{4.6cm}|C{1.7cm}|C{1.7cm}|Y|}{La paleta institucional: veintiocho tokens en los dos modos, con la procedencia de cada valor}{\thd{Token} & \thd{Claro} & \thd{Oscuro} & \thd{Procedencia}}
  \texttt{ground} & \texttt{FFFFFF} & \texttt{0B1B2B} & Claro: superficie del archivo. Oscuro: derivado del color de navegación bajando su luminancia; azul profundo y nunca negro puro \\
  \texttt{ground-alt} & \texttt{F1F4F8} & \texttt{102A43} & Claro: derivado del color de navegación, aclarado hasta 1{,}10:1 sobre la superficie. Oscuro: el color de navegación usado como panel \\
  \texttt{grid} & \texttt{DCE3EC} & \texttt{1D3348} & Derivado en los dos modos: un paso del suelo hacia la rampa. Decorativo por definición \\
  \texttt{reticula} & \texttt{FFFFFF} & \texttt{0B1B2B} & Su propio suelo: bajo este tema la cuadrícula modular no se pinta, y se declara con un color real en lugar de una palabra clave que el cálculo no sabría medir \\
  \texttt{corriente-apagado} & \texttt{A8B4C4} & \texttt{46586B} & Derivado. El peldaño se queda a propósito por debajo de 3:1, porque no informa \\
  \texttt{corriente-tenue} & \texttt{4A5A6E} & \texttt{93A6BC} & Derivado del color de navegación hacia el suelo hasta cruzar 4{,}5:1 \\
  \texttt{corriente-medio} & \texttt{1D4C6E} & \texttt{BFD0E2} & Claro: color secundario del archivo. Oscuro: el mismo matiz elevado sobre el suelo profundo \\
  \texttt{corriente-pleno} & \texttt{102A43} & \texttt{EAF2FA} & Claro: color de navegación del archivo. Oscuro: el mismo matiz elevado \\
  \texttt{accion} & \texttt{086B70} & \texttt{3FB3B5} & Claro: color de acción del archivo, con su valor literal. Oscuro: elevado sobre el suelo profundo hasta cruzar el listón que necesita el rótulo de un botón \\
  \texttt{accion-apoyo} & \texttt{15989A} & \texttt{5FD0D2} & Claro: color de apoyo del archivo, con su valor literal. Oscuro: el mismo matiz elevado \\
  \texttt{seleccion} & \texttt{E6F2F1} & \texttt{123443} & Tinte del suelo hacia la acción. Es superficie y nunca lleva texto propio \\
  \texttt{barra-lateral} & \texttt{102A43} & \texttt{102A43} & Color de navegación en los dos modos: el archivo le da a ese azul exactamente un trabajo, y la barra es la navegación \\
  \texttt{barra-lateral-activo} & \texttt{086B70} & \texttt{3FB3B5} & El color de acción, relleno: el módulo en pantalla es la selección de la barra \\
  \texttt{barra-lateral-texto} & \texttt{BFD0E2} & \texttt{BFD0E2} & Derivado: el peldaño claro de la rampa institucional, medido sobre la barra y no sobre el suelo de la página \\
  \texttt{barra-lateral-activo-texto} & \texttt{FFFFFF} & \texttt{0B1B2B} & Sobre el bloque relleno: superficie en claro y suelo profundo en oscuro, porque el bloque se eleva con el modo \\
  \texttt{error} & \texttt{933632} & \texttt{EC5B51} & Derivado del color de error del archivo, oscurecido conservando matiz y saturación \\
  \texttt{aviso} & \texttt{A36A10} & \texttt{E8A33D} & Claro: color de atención oscurecido un 12 \%. Oscuro: el mismo color aclarado para el suelo profundo \\
  \texttt{ok} & \texttt{287A58} & \texttt{5FCB94} & Claro: color de éxito del archivo, con su valor literal. Oscuro: el mismo color aclarado \\
  \texttt{info} & \texttt{17395B} & \texttt{B9A7F2} & Claro: color secundario oscurecido hasta cruzar 4{,}5:1. Oscuro: violeta declarado fuera del octeto, por la razón que la subsección de divergencias documenta \\
  \texttt{certificacion-certificado} & \texttt{287A58} & \texttt{5FCB94} & Toma prestado el canal de confirmación entero, valor e icono \\
  \texttt{certificacion-en-revision} & \texttt{A36A10} & \texttt{E8A33D} & Toma prestado el canal de aviso entero \\
  \texttt{certificacion-obsoleto} & \texttt{933632} & \texttt{EC5B51} & Toma prestado el canal de error entero \\
  \texttt{serie-1} & \texttt{1D4ED8} & \texttt{7DD3FC} & Heredado del tema de omisión \\
  \texttt{serie-2} & \texttt{B45309} & \texttt{FFC233} & Heredado del tema de omisión \\
  \texttt{serie-3} & \texttt{1F6F43} & \texttt{4ADE80} & Heredado del tema de omisión \\
  \texttt{serie-4} & \texttt{6D28D9} & \texttt{C4B5FD} & Heredado del tema de omisión \\
  \texttt{serie-5} & \texttt{0E7490} & \texttt{67E8F9} & Heredado del tema de omisión \\
  \texttt{serie-6} & \texttt{9D174D} & \texttt{F9A8D4} & Heredado del tema de omisión \\
\end{uxtablalarga}

Los seis colores de serie \textbf{no cambian con el tema}, y es una decisión y no un olvido: una
serie es un canal de datos y no identidad de marca, el archivo de diseño no declara paleta
categórica y las seis ya están verificadas bajo tres dicromacias. Lo que sí se vuelve a medir por
tema es su razón sobre cada suelo, porque el suelo oscuro institucional es \texttt{0B1B2B} y no el
del tema de omisión: las seis quedan por encima de 3:1 en las cuatro combinaciones, con un mínimo
de \textbf{5{,}02:1} en claro y \textbf{9{,}43:1} en oscuro. Contando los dos temas, las
\textbf{veinticuatro} mediciones quedan por encima del umbral, con un mínimo de \textbf{4{,}64:1}.

\subsection{Los tres estados de certificación del catálogo}

Una familia merece nombre propio, porque su regla es de sentido y no de color: las tres etiquetas de
certificación significan cosas opuestas para quien decide si usa un campo, de modo que cada estado
toma prestado un \textbf{canal semántico entero} ---valor, icono y clase---. Así \enquote{tres
estados, tres canales} es estructural y no una coincidencia entre tres códigos que hoy difieren.

\begin{uxtabla}{|L{3.0cm}|L{3.2cm}|L{2.6cm}|Y|}{Los tres estados de certificación del catálogo}
  \uxheadrow \thd{Estado} & \thd{Icono} & \thd{Canal} & \thd{Qué le dice a quien consulta} \\
  Certificado & Marca de verificación & Confirmación & El campo se puede usar \\
  En revisión & Reloj & Aviso & Se puede usar con reserva \\
  Obsoleto & Círculo tachado & Error & \textbf{No se debe usar} \\
\end{uxtabla}

El icono viaja con el color en la definición del token y no se vuelve a elegir en el componente, que
es la regla que impide que dos estados opuestos compartan forma. La consecuencia medible es que las
tres distancias de esta familia son un subconjunto de las seis del canal semántico, de modo que el
piso publicado en la subsección de dicromacia las cubre sin excepción.

\subsection{La familia tipográfica del tema}

El eje del tema son dos cosas: el color y la familia tipográfica. El archivo de diseño declara la
suya con su propio criterio: \enquote{Inter mantiene claridad en tablas, filtros y metadatos. Los
tamaños priorizan densidad sin reducir la legibilidad; el espaciado entre letras permanece en 0.}

Bajo el tema institucional, los roles de titular y de cuerpo se componen con \textbf{Inter},
servida por el mismo módulo de fuentes y con el mismo proveedor que las dos familias del tema de
omisión, que reparte \textbf{Lexend Deca} en titulares y \textbf{Fira Sans} en cuerpo. La familia
monoespaciada \textbf{no cambia}: el rol de dato existe por el ancho fijo de las cifras y no por
identidad, y el archivo de diseño no declara ninguna monoespaciada. Los nueve roles conservan
tamaño, interlínea y espaciado entre letras, de modo que el tema cambia la letra y no la jerarquía.

\figuraux{figuras/a4/guia_tipografia.png}{Lámina de tipografía del archivo de diseño, con el
criterio que declara la familia y su aplicación a títulos, métricas, contenido y controles. Es
diseño de alta fidelidad y no una captura del portal en ejecución.}{fig:a4-tema-tipografia}

\subsection{Matriz de contraste del tema institucional, sobre lo que se lee encima}

El cálculo es el de la subsección 3.3 ---la fórmula de luminancia relativa de la norma contra sus
umbrales de 4{,}5:1 para texto normal, 3:1 para elemento gráfico y 7:1 para el nivel más alto (W3C,
2023)--- aplicado ahora a los veintiocho tokens del tema institucional. Lo que cambia es el
referente: cada token se mide contra \textbf{lo que de verdad se lee encima}, y no siempre es el
suelo de la página. Las etiquetas de la barra lateral se miden sobre la barra, porque bajo este
tema la barra es azul profundo mientras la página es blanca, y una razón tomada sobre la página
sería un número que el lector no ve nunca.

\begin{uxtablalarga}{|L{5.0cm}|C{1.7cm}|C{1.8cm}|C{1.7cm}|Y|}{Matriz de contraste del tema institucional, cada token sobre lo que se lee encima}{\thd{Token} & \thd{Claro} & \thd{Veredicto} & \thd{Oscuro} & \thd{Veredicto}}
  \texttt{ground-alt} & 1{,}10:1 & superficie & 1{,}19:1 & superficie \\
  \texttt{grid} & 1{,}29:1 & decorativo & 1{,}34:1 & decorativo \\
  \texttt{reticula} & 1{,}00:1 & decorativo & 1{,}00:1 & decorativo \\
  \texttt{corriente-apagado} & 2{,}10:1 & decorativo & 2{,}38:1 & decorativo \\
  \texttt{corriente-tenue} & 7{,}05:1 & AAA & 6{,}98:1 & AA \\
  \texttt{corriente-medio} & 9{,}09:1 & AAA & 11{,}06:1 & AAA \\
  \texttt{corriente-pleno} & 14{,}64:1 & AAA & 15{,}41:1 & AAA \\
  \texttt{accion} & 6{,}27:1 & AA & 6{,}90:1 & AA \\
  \texttt{accion-apoyo} & 3{,}51:1 & AA texto grande & 9{,}49:1 & AAA \\
  \texttt{seleccion} & 1{,}15:1 & decorativo & 1{,}33:1 & decorativo \\
  \texttt{barra-lateral} & 14{,}64:1 & superficie & 1{,}19:1 & superficie \\
  \texttt{barra-lateral-activo} & 6{,}27:1 & superficie & 6{,}90:1 & superficie \\
  \texttt{barra-lateral-texto} \textit{(sobre la barra)} & 9{,}30:1 & AAA & 9{,}30:1 & AAA \\
  \texttt{barra-lateral-activo-texto} \textit{(sobre su bloque)} & 6{,}27:1 & AA & 6{,}90:1 & AA \\
  \texttt{error} & 7{,}46:1 & AAA & 5{,}13:1 & AA \\
  \texttt{aviso} & 4{,}54:1 & AA & 8{,}07:1 & AAA \\
  \texttt{ok} & 5{,}23:1 & AA & 8{,}67:1 & AAA \\
  \texttt{info} & 11{,}85:1 & AAA & 8{,}19:1 & AAA \\
  \texttt{certificacion-certificado} & 5{,}23:1 & AA & 8{,}67:1 & AAA \\
  \texttt{certificacion-en-revision} & 4{,}54:1 & AA & 8{,}07:1 & AAA \\
  \texttt{certificacion-obsoleto} & 7{,}46:1 & AAA & 5{,}13:1 & AA \\
  \texttt{serie-1} & 6{,}70:1 & AA & 10{,}44:1 & AAA \\
  \texttt{serie-2} & 5{,}02:1 & AA & 10{,}79:1 & AAA \\
  \texttt{serie-3} & 6{,}15:1 & AA & 9{,}99:1 & AAA \\
  \texttt{serie-4} & 7{,}10:1 & AAA & 9{,}43:1 & AAA \\
  \texttt{serie-5} & 5{,}36:1 & AA & 12{,}01:1 & AAA \\
  \texttt{serie-6} & 7{,}88:1 & AAA & 9{,}60:1 & AAA \\
\end{uxtablalarga}

Tres veredictos de la tabla no son grados de la norma sino papeles del sistema.
\textbf{Superficie} marca un fondo medido contra otro fondo: nada se lee nunca con la fila alterna
ni con la barra lateral como primer plano, así que exigirles 4{,}5:1 sería acusar al sistema de un
incumplimiento que no existe. \textbf{Decorativo} marca los cuatro tokens que declaran no informar
---retícula, filete, conector en reposo y superficie elegida---, y su exención es la razón misma por
la que esa declaración existe en el contrato. \textbf{AA texto grande} marca el único token que
cumple 3:1 y no 4{,}5:1: el verde azulado de apoyo, que el archivo reserva para realce, subrayado y
filete de foco, y que por eso nunca lleva prosa a tamaño de cuerpo.

El resultado agregado es el que importa: \textbf{cero incumplimientos} en las cuatro combinaciones
de tema y modo, sobre los veintiocho tokens. El mínimo real de toda la matriz es \textbf{4{,}54:1},
el del canal de aviso bajo el tema institucional en claro. Ningún token que informa se queda por
debajo de su listón y ningún token que declara no informar lo alcanza, que son las dos formas de
romper la regla.

\subsection{Separación bajo dicromacia, medida por pares}

El contraste sobre el suelo no basta para el canal semántico: dos marcas pueden cumplir el umbral
por separado y confundirse entre sí para quien tiene una dicromacia. Por eso los cuatro semánticos
se miden también dos a dos, simulando protanopia, deuteranopia y tritanopia, y de cada par se
publica la distancia \textbf{peor} de las tres.

\begin{uxtabla}{|L{3.2cm}|C{1.7cm}|L{2.6cm}|C{1.7cm}|Y|}{Separación de los pares semánticos del tema institucional, con la dicromacia que produce la peor distancia}
  \uxheadrow \thd{Par} & \thd{Claro} & \thd{Dicromacia} & \thd{Oscuro} & \thd{Dicromacia} \\
  Error y aviso & 20{,}7 & tritanopia & 22{,}6 & deuteranopia \\
  Error y confirmación & 14{,}5 & protanopia & 23{,}3 & protanopia \\
  Error e informativo & 45{,}4 & protanopia & 67{,}6 & tritanopia \\
  Aviso y confirmación & 37{,}8 & protanopia & 37{,}8 & protanopia \\
  Aviso e informativo & 54{,}3 & tritanopia & 35{,}3 & tritanopia \\
  Confirmación e informativo & 22{,}8 & tritanopia & 25{,}9 & tritanopia \\
\end{uxtabla}

El piso no es un número inventado para esta paleta: es la peor separación que el tema de omisión ya
sostiene, \textbf{13{,}6} en claro y \textbf{21{,}5} en oscuro. El tema institucional queda por
encima de los dos, con \textbf{14{,}5} y \textbf{22{,}6}. Un tema nuevo que empeorara ese número
habría degradado el sistema aunque cada token cumpliera su umbral por separado.

Los tres estados de certificación se miden como familia aparte, porque la distancia solo es una
exigencia \emph{dentro} de una familia: un estado del catálogo nunca comparte superficie con un
mensaje de error del formulario. Sus tres distancias son \textbf{14{,}5} en claro y \textbf{22{,}6}
en oscuro, idénticas a las del canal semántico por la razón estructural ya dicha ---toman prestados
esos mismos canales--- y por encima del mismo piso.

\subsection{Las dos correcciones de luminancia sobre el octeto}

Dos colores del octeto entran al contrato con la luminancia corregida, y la corrección vive en el
valor del token y no en la prosa. La norma del sistema es que \textbf{ningún valor derivado se
publica con el rótulo de literal}: un color corregido que se presenta como oficial es el defecto
caro, porque nadie lo vuelve a revisar.

\begin{enumerate}[leftmargin=*,itemsep=5pt,topsep=6pt]
\item \textbf{El ámbar de atención no alcanza el listón de texto con su valor literal.} El valor del archivo,
  \texttt{B97812}, da \textbf{3{,}65:1} sobre la superficie blanca, por debajo del 4{,}5:1 que exige
  un token que informa. Oscurecido un \textbf{12 \%} conservando matiz y saturación llega a
  \texttt{A36A10} con \textbf{4{,}54:1}. Es el mínimo que cruza: a un 10 \% todavía da 4{,}40:1.
  Publicar el valor del archivo exigiría un tercer nivel de contrato ---uno para componente y texto
  grande--- y esa decisión de sistema queda declarada como pendiente, no resuelta a escondidas.
\item \textbf{El rojo de error no se embarca con su valor literal.} Con \texttt{B8443F} el error da
  5{,}33:1 y la confirmación 5{,}23:1, de modo que bajo protanopia los dos se separan por
  \textbf{10{,}1} contra un piso de 13{,}6: un error se leería como una confirmación. Oscurecido un
  \textbf{20 \%}, \texttt{933632} sube a 7{,}46:1 y su peor par a \textbf{14{,}5}. En oscuro, el
  mismo color aclarado para el suelo profundo se queda en \textbf{13{,}1} contra un piso de
  21{,}5; oscurecido un 12 \%, \texttt{EC5B51} da 5{,}13:1 y su peor par sube a \textbf{22{,}6}.
\end{enumerate}

Hay además una derivación que no es una corrección del archivo sino una ausencia suya: el
\textbf{canal informativo}. En modo claro sale del octeto, del color secundario \texttt{1D4C6E}
oscurecido hasta cruzar 4{,}5:1, y su distancia a la confirmación queda en 22{,}8. En modo oscuro
ningún azul sirve, porque sobre un suelo profundo el azul y el verde convergen: el par informativo y
confirmación cae a 10{,}4 bajo tritanopia contra un piso de 21{,}5. El violeta \texttt{B9A7F2} es el
matiz más cercano que sobrevive, con 25{,}9, y se declara como salida del octeto en lugar de
presentarse como uno de sus colores.

Las tres derivaciones siguen la misma regla, escrita antes de aplicarla: se corrige la luminancia y
se conservan el matiz y la saturación, de manera que el color siga siendo reconocible como el del
archivo. Y las tres se descubrieron por cálculo y no a ojo, que es la única forma de encontrar un
defecto que solo existe para un lector con dicromacia.

\subsection{Los ocho estados interactivos, declarados como derivación}

El archivo de diseño publica ocho estados interactivos y los declara por \textbf{nombre de
variable}, no por valor: la lámina muestra el color y nombra la variable que lo produce. Esta
entrega no inventa los ocho valores ni los copia de la lámina con un cuentagotas: los deriva con
una regla escrita y marca la derivación como tal.

\textbf{La regla.} Sobre el color base, un \textbf{12 \%} de oscurecimiento para el estado al pasar
el puntero y un \textbf{20 \%} para el estado presionado, conservando matiz y saturación. Es la
misma operación con la que se corrigieron las divergencias de la subsección anterior, de modo que
el sistema tiene una sola forma de mover un color.

\begin{uxtabla}{|L{4.2cm}|L{3.4cm}|Y|}{Los ocho estados interactivos del archivo de diseño y su procedencia}
  \uxheadrow \thd{Estado} & \thd{Base} & \thd{Procedencia} \\
  Primario al pasar & \texttt{accion} & Derivación: base oscurecida un 12 \% \\
  Primario presionado & \texttt{accion} & Derivación: base oscurecida un 20 \% \\
  Secundario al pasar & \texttt{corriente-pleno} & Derivación: base oscurecida un 12 \% \\
  Secundario presionado & \texttt{corriente-pleno} & Derivación: base oscurecida un 20 \% \\
  Deshabilitado & \texttt{corriente-apagado} & Derivación: peldaño apagado, sin papel informativo \\
  Foco & \texttt{accion-apoyo} & Variable propia del archivo, aplicada como filete de foco \\
  Error & \texttt{error} & Variable propia del archivo, con el valor corregido del token \\
  Información & \texttt{info} & Variable propia del archivo, con el valor derivado del token \\
\end{uxtabla}

\figuraux{figuras/a4/guia_paleta.png}{Lámina de paleta y estados interactivos del archivo de
diseño. Los ocho estados se declaran por nombre de variable, que es la razón por la que esta
entrega los publica como derivación con su regla. Es diseño de alta fidelidad y no una captura del
portal en ejecución.}{fig:a4-tema-paleta}

Las cuatro primeras filas no dicen \enquote{color de acción} en abstracto: dicen el nombre del token
que las produce, porque ese color tiene ranura propia en el contrato. La norma es que \textbf{un
estado interactivo se deriva de un token y nunca de un valor escrito en la capa de componentes}, que
es la vía por la que un color se duplica a mano en tres sitios y se separa de su origen sin que nadie
lo note.

\subsection{El tema, medido en el portal en ejecución}

Las cifras de las tablas anteriores salen del generador. Lo que sigue sale del portal abierto en un
navegador sin cabeza, que es la única forma de comprobar que lo calculado es también lo pintado.
Sobre las cuatro combinaciones se leyó el valor resuelto de cada propiedad personalizada en la raíz
del documento, la familia tipográfica computada del cuerpo y del titular, y el número de elementos
cuyo contenido desborda su caja.

\begin{uxtabla}{|L{3.0cm}|C{1.9cm}|C{1.9cm}|C{2.0cm}|C{1.9cm}|Y|}{Medición en el navegador de las cuatro combinaciones, a 1440 por 900 píxeles}
  \uxheadrow \thd{Combinación} & \thd{Suelo} & \thd{Barra lateral} & \thd{Acción} & \thd{Familia} & \thd{Desbordes} \\
  Omisión, claro & \texttt{F4F6F9} & \texttt{EAEEF4} & \texttt{14171D} & Fira Sans & 0 \\
  Omisión, oscuro & \texttt{0A0A0C} & \texttt{131519} & \texttt{E8F4FF} & Fira Sans & 0 \\
  Institucional, claro & \texttt{FFFFFF} & \texttt{102A43} & \texttt{086B70} & Inter & 0 \\
  Institucional, oscuro & \texttt{0B1B2B} & \texttt{102A43} & \texttt{3FB3B5} & Inter & 0 \\
\end{uxtabla}

Las cuatro pintan su propia paleta y ninguna hereda la del tema vecino. Los cuatro valores de la
columna del suelo y los cuatro de la columna de acción coinciden exactamente con los que la tabla de
procedencia declara, que es la comprobación que importa: entre la definición y el píxel hay un
generador, una hoja de estilo y un servidor, y cualquiera de los tres podría haber perdido un valor
en el camino. La familia cambia con el tema y no con el modo, que es lo que el eje declara, y la
retícula la absorbe sin un solo corte: Inter tiene mayor altura de x que Fira Sans y el riesgo real
era que una celda densa creciera.

En la misma corrida se midió el \textbf{desplazamiento horizontal} en las diez rutas del chasis, a
390, 768, 1280 y 1440 píxeles de ancho, y en las cuatro combinaciones de tema y modo: \textbf{160
mediciones con cero exceso}. La sección de validación publica esa medición con su línea base y su
regla.

La preferencia se comprobó además contra el HTML que devuelve el servidor, y no contra el que existe
después de hidratar: con las dos cookies puestas la etiqueta de apertura llega ya como
\texttt{<html lang=es data-tema=institucional data-modo=oscuro>}, con sus tres valores entre
comillas. Es la diferencia entre una
preferencia que se aplica y una que se corrige a la vista del lector. El tema de omisión no escribe
ese atributo, y tampoco es un descuido: ese tema \emph{es} el bloque base de la hoja generada y no
una anulación de nada, de modo que la ausencia del atributo es su forma correcta de declararse.

\subsection{El eje de tema y modo, sobre una misma pantalla}

Las cuatro imágenes siguientes recorren el eje completo ---dos temas por dos modos--- sobre
\textbf{la misma pantalla, la misma sesión y el mismo contenido}, que es la única disposición en la
que lo comparado es el sistema consigo mismo. Se tomaron en una sola pasada guionizada contra el
portal en ejecución, a 1440 por 900 píxeles y en español. El guion comprueba antes de disparar que
la raíz del documento declara el tema y el modo que la imagen dice mostrar, de modo que una cookie
ignorada produce un fallo y no una captura plausible. Las capturas del tema institucional pantalla
por pantalla se publican en la sección de prototipos.

Las cuatro comparten contenido, estructura y orden de tabulación. Lo único que cambia entre ellas es
el par tema y modo, que es exactamente lo que el eje declara: si algo más cambiara, el eje estaría
haciendo un trabajo que no le corresponde. Las dos primeras y las dos últimas son además la
comparación que sostiene la razón de ser de cada tema: la misma información, separada por luminancia
en un caso y por identidad cromática en el otro.

\figuraux[0.78\linewidth]{figuras/a4/tema/combinacion_corriente_claro_inicio.png}{Espacio de trabajo
operativo bajo el \textbf{tema de omisión en modo claro}. La cuadrícula modular está pintada y el
módulo en curso de la barra lateral se distingue por luminancia y peso, sin bloque relleno. Captura
del portal en ejecución.}{fig:a4-tema-corriente-claro}

\figuraux[0.78\linewidth]{figuras/a4/tema/combinacion_corriente_oscuro_inicio.png}{La misma pantalla
bajo el \textbf{tema de omisión en modo oscuro}, sobre el suelo \texttt{0A0A0C}. Cambia el modo y no
el tema: la retícula sigue pintada y la familia sigue siendo Fira Sans. Captura del portal en
ejecución.}{fig:a4-tema-corriente-oscuro}

\figuraux[0.78\linewidth]{figuras/a4/tema/combinacion_institucional_claro_inicio.png}{La misma
pantalla bajo el \textbf{tema institucional en modo claro}. La barra lateral pasa al azul profundo
del archivo, el módulo en curso es un bloque relleno de verde azulado y la retícula desaparece.
Captura del portal en ejecución.}{fig:a4-tema-inst-claro}

\figuraux[0.78\linewidth]{figuras/a4/tema/combinacion_institucional_oscuro_inicio.png}{La misma
pantalla bajo el \textbf{tema institucional en modo oscuro}, sobre el suelo \texttt{0B1B2B}. El azul
profundo sustituye al negro puro por la regla que el sistema ya se había puesto. Captura del portal
en ejecución.}{fig:a4-tema-inst-oscuro}

\clearpage

% ============================================================================
\section{Componentes gráficos de la interface}

Los componentes de esta sección son los que la ruta \texttt{/guia} renderiza: botones, campos con sus chips e insignias, tablas y tarjetas. La ruta es el sistema de diseño vivo del producto, no una maqueta del sistema: cada lámina lee la paleta tipada que el generador emite, así que un cambio de token se ve en la lámina sin tocar el componente. El armazón de navegación ---barra lateral de 240 px con dos niveles, cabecera de 56 px y franja de alcance--- se documenta con las capturas de las siete pantallas, porque es donde se usa; aquí quedan sus medidas, en la sección 7.

El sistema fija además una postura sobre la personalización: los valores por omisión los decide el perfil y no el usuario. El respaldo es empírico y ajeno al equipo: en un estudio con 20 diseñadores que juzgaron por pares 600 interfaces generadas, el acuerdo entre ellos fue de $\kappa = 0{,}25$ sobre dimensiones tan básicas como densidad de información y jerarquía (Peng et al., 2026). Con esa dispersión, dejar la apariencia a criterio individual produce tantas interfaces como usuarios; el sistema prefiere un valor por omisión defendible por rol.

\subsection{Botones: matriz de 17 celdas}

Tres variantes por cinco estados, más dos celdas adicionales: \textbf{15 + 2 = 17}, frente a las nueve del material de referencia. La lámina de botones de \texttt{/guia} las renderiza todas a la vez, de modo que una regresión de estado se ve sin recorrer la aplicación.

\begin{uxtabla}{|L{1.7cm}|Y|Y|Y|Y|Y|}{Matriz de botones: tres variantes por cinco estados}
  \uxheadrow \thd{Variante} & \thd{Normal} & \thd{Puntero encima} & \thd{Foco visible} & \thd{Activo} & \thd{Deshabilitado} \\
  Contenido & Relleno primaria 500 y texto de superficie, 4{,}94:1 & Relleno primaria 700 & Anillo de 2 px separado 2 px del borde & Relleno primaria 900 & Relleno de filete y texto atenuado, sin puntero \\
  Contorno & Borde \texttt{line-strong} y texto primaria 500 & Fondo de superficie alterna & Mismo anillo de 2 px & Fondo primaria 100 & Borde de filete y texto atenuado \\
  Texto & Solo texto primaria 500 & Subrayado & Mismo anillo de 2 px & Texto primaria 700 & Texto atenuado \\
\end{uxtabla}

\begin{uxtabla}{|L{2.8cm}|Y|}{Las dos celdas adicionales de la matriz de botones}
  \uxheadrow \thd{Celda} & \thd{Tratamiento} \\
  Destructiva & Relleno semántico de error con texto de superficie, 6{,}18:1. Separada espacialmente del grupo de acciones normales y siempre con confirmación previa \\
  En carga & Indicador de progreso a la izquierda del rótulo, botón inhabilitado y \texttt{aria-busy} mientras dura la operación. El rótulo no cambia de texto, para que el ancho del botón no salte \\
\end{uxtabla}

\textbf{Reglas de la lámina.} Altura de 32 px en densidad alta y de 40 px en densidad normal · área táctil mínima de 44 × 44 px aunque el botón se vea más pequeño · el rótulo nunca se parte en dos líneas · \textbf{una sola acción primaria por pantalla} · icono siempre a la izquierda del texto y alineado a la línea base.

\subsection{Campos de formulario: seis estados}

\begin{uxtabla}{|L{2.4cm}|Y|}{Los seis estados del campo de formulario}
  \uxheadrow \thd{Estado} & \thd{Tratamiento} \\
  Normal & Borde \texttt{line-strong} sobre superficie, 4{,}55:1, con etiqueta visible encima del campo \\
  Con foco & Anillo de 2 px en primaria 500 y borde del mismo color; el anillo no sustituye al borde, se suma \\
  Con valor & Igual que el normal, con el texto en \texttt{ink}; el marcador de posición desaparece y la etiqueta permanece \\
  Error & Borde y texto de ayuda en el semántico de error, mensaje debajo del campo y anuncio con \texttt{role="alert"} \\
  Deshabilitado & Fondo de superficie alterna, texto atenuado, sin puntero y fuera del orden de tabulación \\
  Solo lectura & Sin borde de acción y con fondo de superficie alterna, pero \textbf{seleccionable y dentro del orden de tabulación}: se ve distinto del deshabilitado porque significa algo distinto \\
\end{uxtabla}

\textbf{Reglas de la lámina.} Etiqueta visible sobre el campo, nunca solo el marcador de posición · texto de ayuda persistente bajo los campos complejos · el error va debajo del campo que falló y no en un resumen al principio · la validación ocurre al salir del campo, no en cada tecla. El borde de \texttt{line-strong} no es una preferencia estética: es el hallazgo 4 de la sección 3.4 convertido en componente.

\subsection{Chips e insignias}

Cinco semánticas de chip ---neutro, informativo, éxito, aviso y error---, las cuatro insignias de rol ---operativo, analista, directivo y administrador--- y el chip de faceta transversal, que lleva su propio icono. \textbf{Todos llevan icono además de color}, por la regla de no depender del color que la sección 9 fija. El chip de aviso usa el acento 900 sobre el acento 100, que es la regla sustituta del hallazgo 2, y las insignias de rol usan el radio completo, que es la única excepción declarada de la escala de radios.

\subsection{Tablas de datos}

\begin{uxtabla}{|L{4.2cm}|Y|}{Reglas de la tabla de datos}
  \uxheadrow \thd{Regla} & \thd{Detalle} \\
  Cabecera fija & Permanece visible al desplazar el cuerpo de la tabla \\
  Altura de fila & \texttt{--table-row-height}, 36 px, que es la densidad del producto \\
  Fila alterna & Superficie alterna; el texto de cuerpo sobre ella mantiene 13{,}12:1 \\
  Ordenamiento & \texttt{aria-sort} refleja el estado real de la columna, no solo el icono \\
  Cifras & Alineadas a la derecha y con cifras tabulares \\
  Acciones & Columna anclada a la derecha, fuera del área de lectura de datos \\
  Volumen & Paginación o virtualización a partir de 50 filas \\
\end{uxtabla}

\subsection{Tarjetas: KPI y tarjeta de llamada a herramienta}

La \textbf{tarjeta KPI} lleva rótulo, cifra grande en cifras tabulares, variación con signo, icono de dirección y micrográfica de tendencia. La variación nunca se comunica solo con color: lleva flecha y signo.

La \textbf{tarjeta de llamada a herramienta} es la pieza que sostiene la regla anti-alucinación del producto ---toda cifra que aparece en una respuesta procede de una consulta visible y cita su fuente--- y la lámina de tarjetas la muestra en sus cuatro estados como galería. La decisión de hacer visible cada llamada tiene respaldo publicado: en un estudio con 13 personas que construyen y usan sistemas multiagente, la transparencia aparece definida de cinco maneras distintas ---reproducibilidad, depuración, delimitación, visualización y auditoría--- y las cinco exigen ver qué hizo el sistema, no solo qué respondió (Naik et al., 2026). El mismo argumento sostiene el estado compartido entre el tablero y el asistente, donde mantener una sola representación del estado analítico redujo la tasa de agotamiento de tiempo del 40{,}0 \% al 10{,}0 \% en el experimento controlado de TwinBI (Jang y Li, 2026).

\begin{uxtabla}{|L{2.4cm}|L{4.4cm}|Y|}{Los cuatro estados de la tarjeta de llamada a herramienta}
  \uxheadrow \thd{Estado} & \thd{Qué muestra} & \thd{Regla} \\
  Anuncio & Nombre legible de la operación y qué fuente consulta & Aparece antes de tener el dato: el usuario sabe qué se está preguntando \\
  Ejecución & Indicador de progreso y tiempo transcurrido & El tiempo es real y se actualiza; no hay animación que finja avance \\
  Resultado & Cifra o mini tabla plegable con la cita de la fuente del catálogo & El resultado se muestra \textbf{antes} del texto generado \\
  Error & Qué paso falló y qué hacer & Nombra la causa y ofrece la acción de reintento; no dice solo que hubo un error \\
\end{uxtabla}

\subsection{Anatomía compartida}

Las once láminas comparten forma: cada una expone un atributo de datos propio para que el guion de captura la recorte sin adivinar, y ninguna contiene lógica de negocio. Presentan lo que el generador emitió. Esa simetría es lo que permite que la guía se capture con un guion reproducible en lugar de con recortes a mano de distinto tamaño.

% ============================================================================
\section{Iconografía}

\subsection{Una sola familia}

\textbf{Lucide, para todo el producto.} No se mezclan familias y no se dibujan iconos a mano. La razón conviene escribirla porque no es obvia: mezclar familias produce grosores de trazo y radios de esquina inconsistentes que el ojo detecta aunque el lector no sepa nombrar por qué algo se ve poco cuidado. La familia se instala como paquete y los iconos viajan en el propio paquete de la aplicación, de modo que la interfaz no depende de un servicio externo para dibujarse.

\subsection{Tokens de tamaño y grosor}

Tres tamaños y ninguno intermedio: \textbf{16, 20 y 24 px}. Grosor de trazo de \textbf{1{,}5 px en los tres}, sin engrosar en el tamaño mayor. Alineación a la línea base del texto que acompañan, no al centro óptico de la caja.

\subsection{Contorno y relleno}

Contorno en toda la interfaz. El relleno queda reservado en exclusiva al \textbf{estado activo de la navegación}, que es la única distinción de forma que el producto se permite; en cualquier otro lugar, un icono relleno significaría algo que el sistema no define.

\subsection{Área táctil y rótulo accesible}

Área táctil de 44 × 44 px aunque el icono mida 16. \textbf{Rótulo accesible obligatorio en todo icono sin texto}, y el icono que acompaña a un texto se marca como decorativo para que el lector de pantalla no lo lea dos veces. Contraste mínimo de 3:1 contra su fondo, que es el mismo umbral con el que la sección 3 juzga cualquier elemento gráfico. \textbf{Nunca un emoji como icono estructural}: su forma cambia con el sistema operativo y su lectura en voz alta es impredecible.

\subsection{Inventario por función}

\begin{uxtabla}{|L{2.8cm}|Y|}{Inventario de iconos del producto, agrupado por función}
  \uxheadrow \thd{Función} & \thd{Iconos} \\
  Navegación & Los cuatro módulos del mapa de A3: inicio, exploración y extracción, gobierno del dato y administración \\
  Acciones & Buscar, filtrar, exportar, compartir y copiar \\
  Estado & Vigente, en revisión, retirado y error \\
  Datos & Tabla, gráfica, linaje y fuente \\
  Gobierno & Diccionario, calidad y carga \\
  Administración & Usuario, rol, llave y bitácora \\
  Asistente & Enviar, detener, herramienta y reintentar \\
  Interfaz & Marcador de rama del árbol de navegación y selector de idioma, los dos ya en uso en la aplicación \\
\end{uxtabla}

\subsection{Cómo se declara el inventario}

El inventario se declara \textbf{una sola vez}, como lista en la configuración del módulo de iconos, y la lámina lo recorre desde esa misma lista. El motivo es concreto y costó descubrirlo: el empaquetado automático detecta los nombres escritos literalmente en las plantillas, pero no los que se arman en tiempo de ejecución al concatenar el prefijo de la familia con el nombre del icono. Con la lista declarada, la lámina de iconografía dibuja los mismos iconos en desarrollo y en producción; sin ella, la versión construida muestra huecos donde el navegador esperaba un trazo.

% ============================================================================
\section{Imágenes e ilustraciones}

\subsection{Regla de origen}

\textbf{Los modelos de imagen se usan solo para fotografía y marca, nunca para diagramas con texto en español.} La regla costó descubrirla en la Actividad 1 y aquí queda como norma del sistema: las figuras con etiquetas se generan con código ---matplotlib para gráficas de datos, HTML y CSS para retículas, LaTeX para las láminas de esta guía--- y los retratos y escenas de contexto con modelo de imagen, con lista negativa explícita de texto, letras, números y logotipos.

\subsection{Qué imágenes existen en la interfaz}

Tres, y ninguna más: la marca, los avatares construidos a partir de las iniciales del usuario y los gráficos de datos. \textbf{No hay fotografía decorativa en el producto.} Una fotografía de archivo en un portal de datos financieros ocupa el lugar de un dato y no aporta ninguno.

\subsection{Inventario de imágenes del documento}

Las imágenes de la serie de entregables se mantienen versionadas en la carpeta \texttt{imagenes/} de los entregables y suman dieciséis archivos: los ocho retratos de las personas de la Actividad 1, las seis escenas de contexto de los escenarios de la Actividad 2 y los dos logotipos de la portada. Todas se insertan en el documento a través de macros de estilo, nunca con medidas escritas caso por caso.

\subsection{Texto alternativo}

Toda imagen significativa lleva texto alternativo; las decorativas se marcan como tales para que el lector de pantalla las omita. En el documento, cada figura lleva pie numerado, y en la interfaz cada gráfica cumple además la regla de la sección 9: alternativa en tabla y resumen textual.

% ============================================================================
\section{Retícula, espaciado y diseño adaptativo}

\subsection{Retícula de doce columnas}

Doce columnas con un canal de \texttt{--grid-gap} y márgenes de \texttt{--card-padding}. La barra lateral queda fuera de la retícula: ocupa su propio ancho fijo y el contenido empieza después.

\subsection{Ritmo de espaciado, radios, sombras y quiebres}

\laminaReticula

\subsection{La única excepción de la escala de radios}

\texttt{--radius-full} es la excepción declarada: se aplica a las insignias de rol, a los avatares y al punto de estado, y a nada más. Un radio completo en una tarjeta o en un campo rompería la lectura de la retícula, que descansa en esquinas de curvatura corta.

\subsection{Tres niveles de sombra y ni uno más}

Reposo, elevado y flotante. Mezclar sistemas de elevación es el error que más rápido delata una interfaz sin sistema: dos sombras parecidas pero distintas se leen como un descuido, no como una jerarquía. Los tres niveles se calculan sobre el mismo neutro oscuro con opacidades crecientes, de modo que la familia se mantiene.

\subsection{Comportamiento adaptativo}

El producto es de escritorio y la retícula lo asume. Por debajo de 1024 px la barra lateral colapsa a iconos; por debajo de 768 px las tablas pasan a tarjetas apiladas y los campos suben a 16 px. El quiebre de 1440 px es el ancho de captura de las figuras de esta actividad, elegido para que las capturas muestren el producto en su densidad de diseño.

\subsection{Sin desplazamiento horizontal de página}

Nunca. Una tabla ancha se desplaza \textbf{dentro de su propio contenedor}, con la cabecera fija y la columna de acciones anclada. El desplazamiento horizontal de la página entera esconde controles fuera del campo visual y rompe el orden de tabulación que la sección 9 exige.

% ============================================================================
\section{Microinteracciones y animaciones}

\subsection{Duraciones y curvas}

Entre \textbf{150 y 300 ms} para microinteracciones y hasta 400 ms para transiciones complejas; ninguna pasa de 500 ms. Curva \texttt{ease-out} al entrar y \texttt{ease-in} al salir, con la salida al 60--70 \% de la duración de la entrada, porque una salida lenta se percibe como lentitud del sistema y no como elegancia.

\subsection{Solo se animan dos propiedades}

\texttt{transform} y \texttt{opacity}. Nunca \texttt{width}, \texttt{height}, \texttt{top} ni \texttt{left}: obligan a recalcular la maquetación y producen saltos visibles justo cuando la interfaz quiere transmitir control.

\subsection{Inventario de animaciones}

El inventario es corto a propósito y cada entrada declara qué comunica. Una animación que no comunica nada se elimina.

\begin{uxtabla}{|L{5.0cm}|L{2.4cm}|Y|}{Inventario de animaciones del producto}
  \uxheadrow \thd{Animación} & \thd{Duración} & \thd{Qué comunica} \\
  Despliegue del segundo nivel de la barra lateral & 200 ms & Que ese contenido pertenece al módulo abierto \\
  Aparición de la tarjeta de llamada a herramienta & 180 ms & Que el sistema empezó a trabajar \\
  Llegada de fragmento en el flujo de respuesta & Continua & Progreso real, no simulado \\
  Apertura del panel de linaje & 240 ms & De dónde viene el panel y adónde vuelve al cerrarse \\
  Cambio de nivel en el desglose & 220 ms & Que se está bajando dentro de la misma jerarquía \\
  Aviso emergente & 150 ms al entrar y 100 ms al salir & Confirmación sin interrumpir la lectura \\
\end{uxtabla}

\subsection{Movimiento reducido}

\texttt{prefers-reduced-motion} se respeta en las seis: las transiciones se reducen a un cambio de opacidad o desaparecen, y el flujo de respuesta pasa a mostrar el texto completo al terminar en lugar de fragmento a fragmento. La preferencia del sistema operativo gana siempre sobre la animación por omisión.

% ============================================================================
\section{Directrices de accesibilidad}

\subsection{Contraste}

4{,}5:1 en texto y 3:1 en elementos gráficos y de datos, \textbf{verificado por cálculo y publicado como matriz} en la sección 3, según los criterios 1.4.3 y 1.4.11 de WCAG 2.2 (W3C, 2023). El neutro \texttt{muted} queda restringido a texto de 14 px o más, porque su 4{,}55:1 pasa AA por un margen que no admite tamaños menores.

\subsection{Foco visible}

Anillo de 2 px en todo elemento interactivo, con separación de 2 px del borde para que se distinga sobre rellenos oscuros. \textbf{No se elimina nunca.} Los componentes del prototipo lo declaran con la variante \texttt{focus-visible}, de modo que aparece con teclado y no estorba con puntero.

\subsection{Recorrido de teclado}

Orden de tabulación igual al orden visual, sin saltos ni trampas de foco. La aplicación abre con un enlace de salto al contenido principal, presente en los dos idiomas del producto. Los elementos deshabilitados salen del recorrido; los de solo lectura permanecen en él.

\subsection{Estructura y semántica}

Jerarquía de encabezados secuencial, sin saltar niveles. La navegación principal y el acceso transversal llevan rótulo accesible propio, y el elemento activo del árbol es exactamente uno: la rama que comparte ruta con su módulo cede la marca al módulo, para que nunca haya dos elementos anunciados como actuales.

\subsection{La información nunca depende solo del color}

Se acompaña siempre de icono, forma, patrón o texto. La regla gobierna los chips, las insignias, la variación de la tarjeta KPI y las seis series de gráfica, que llevan marcador y patrón de línea propios además de color.

\subsection{Gráficas}

Toda gráfica lleva \textbf{alternativa en tabla} y \textbf{resumen textual}. Una gráfica sin alternativa es un dato que existe solo para quien puede verlo, y en un portal de gobierno del dato eso contradice el propósito del producto.

\subsection{Formularios}

El mensaje de error va debajo del campo que falló, se anuncia con \texttt{role="alert"} y el foco viaja al primer campo inválido. El mensaje dice causa y remedio, según la regla de la sección 10.

\subsection{Escala y área táctil}

El contenido se mantiene legible al 200 \% de escala de texto, sin truncamientos ni superposiciones. Área táctil mínima de 44 × 44 px en todo destino interactivo, incluidos los iconos de 16 px.

\subsection{Idioma declarado}

El documento declara su idioma en el elemento raíz y lo cambia con el idioma de la interfaz, de modo que el lector de pantalla usa la voz y las reglas de pronunciación correctas al pasar de español a inglés.

% ============================================================================
\section{Voz y tono}

\subsection{Principios}

Español neutro. \textbf{Sin emojis en ninguna cadena de la interfaz.} Segunda persona en las acciones: \enquote{Exportar}, no \enquote{Exportación}. Los mensajes de error dicen \textbf{causa y remedio}. Las cifras se acompañan siempre de su fuente. Las proyecciones llevan su etiqueta de método a la vista. No hay superlativos ni lenguaje comercial: es una herramienta de trabajo y el tono que la marca reclama es confiable, preciso y sobrio.

\subsection{Se dice y no se dice}

\begin{uxtabla}{|L{5.6cm}|Y|}{Ocho pares de la guía de redacción de la interfaz}
  \uxheadrow \thd{Se dice} & \thd{No se dice} \\
  Tu perfil no tiene permiso para abrir esta pantalla & Acceso denegado \\
  La dirección solicitada no corresponde a ninguna pantalla del prototipo & Error 404 \\
  El servidor no pudo construir la pantalla & Ha ocurrido un error inesperado \\
  Exportar & Exportación de datos \\
  Prototipo con datos sintéticos, no conectado a sistemas reales & Ambiente de demostración \\
  Proyección lineal sobre datos sintéticos & Predicción \\
  Consultando la base de liquidez & Procesando \\
  Volver al índice de prototipos & Regresar \\
\end{uxtabla}

\subsection{El producto habla en dos lenguas}

La interfaz es bilingüe español e inglés con internacionalización real: los dos catálogos declaran \textbf{79 claves cada uno}, con paridad exacta verificada por prueba automática, y el idioma viaja en una cookie propia sin cambiar ninguna dirección. El arranque es determinista en español y el cambio lo hace la persona, no el navegador. \textbf{El PDF de esta actividad es solo en español}: la decisión bilingüe alcanza a la aplicación, no a los entregables del curso.

\subsection{Traducción de producto, no traducción literal}

El inglés del producto se escribe como interfaz en inglés, no como calco del español. Los nombres propios no se traducen.

\begin{uxtabla}{|L{5.2cm}|Y|}{Equivalencias de la interfaz en sus dos idiomas}
  \uxheadrow \thd{Español} & \thd{Inglés} \\
  Exploración y extracción & \textit{Explore and extract} \\
  Gobierno del dato & \textit{Data governance} \\
  Acceso & \textit{Sign in} \\
  Bitácora de accesos & \textit{Access log} \\
  Saltar al contenido & \textit{Skip to content} \\
  Facetas transversales & \textit{Cross-cutting facets} \\
  Karisma Data & Karisma Data \\
\end{uxtabla}

\subsection{El rótulo de idioma se escribe en su propio idioma}

El selector nombra cada opción con su endónimo ---\enquote{Español} y \enquote{English}--- y marca cada botón con su atributo de idioma. La razón es de accesibilidad antes que de estilo: quien no puede leer la pantalla actual necesita reconocer la salida sin traducirla mentalmente.

% ============================================================================
\section{Documentación y control de versiones}

\subsection{Versión vigente}

\textbf{\versionguia\ · \fechaguia.} La cadena de versión aparece en tres sitios y coincide en los tres: la cabecera del bloque \texttt{@theme} de la aplicación, el manifiesto de tokens y esta sección del documento. Un guion de verificación compara las tres cadenas y falla si alguna se desvía, así que la versión impresa aquí es la que el producto ejecuta.

\subsection{Dónde vive la guía}

En dos lugares que no divergen porque salen del mismo generador: el bloque \texttt{@theme} de la hoja generada \texttt{main.css}, que la aplicación consume, y estas once secciones. Todo cambio entra \textbf{por modificación del generador}, nunca editando el CSS ni el documento por separado. Los dos archivos generados llevan cabecera de advertencia y una comprobación de idempotencia detecta cualquier edición a mano: el generador vuelve a emitir las cuatro salidas y compara contra el disco.

\subsection{Bitácora de cambios}

\begin{uxtabla}{|L{2.3cm}|L{1.3cm}|Y|Y|}{Bitácora de cambios de la guía de estilos}
  \uxheadrow \thd{Fecha} & \thd{Versión} & \thd{Qué cambió} & \thd{Por qué} \\
  26 de julio de 2026 & 0.1 & Se anclan once colores y dos familias tipográficas en la hoja de estilo del curso & La Actividad 1 necesitaba una identidad visual estable para el primer entregable \\
  2 y 9 de agosto de 2026 & 0.1 & Sin cambio de valores & Las Actividades 2 y 3 compilan contra las mismas once anclas; la continuidad visual de la serie es el argumento \\
  10 de agosto de 2026 & 0.9 & La hoja pasa de 11 a 37 nombres de color, de forma aditiva pura: veinte declaraciones nuevas y seis reúsos explícitos & La guía necesita familias completas, neutros, semánticos y serie categórica; renombrar o mover un valor invalidaría tres entregas cerradas \\
  10 de agosto de 2026 & 0.9 & Un generador único emite el bloque \texttt{@theme}, la paleta tipada, las láminas del documento y el manifiesto & Elimina la posibilidad de que la aplicación y el PDF muestren colores distintos \\
  10 de agosto de 2026 & 0.9 & Se calculan 44 pares de contraste y se corrigen los cuatro defectos de la sección 3.4 en los tokens & El cálculo refutó cuatro reglas que la planeación daba por buenas \\
  10 de agosto de 2026 & 0.9 & La interfaz pasa a bilingüe español e inglés y el contrato de navegación deja de guardar palabras & Sin el cambio, el índice de prototipos habría quedado en español dentro de la interfaz en inglés \\
  16 de agosto de 2026 & 1.0 & Publicación de la guía con la Actividad 4 & Primera versión completa: once secciones, once láminas vivas y matriz de contraste publicada \\
\end{uxtabla}

Cada entrada la aprueba el Equipo 8 (Alexandro Mayoral, Jacqueline Sarmiento y Arthur Zizumbo) en la revisión de la entrega semanal, que es el único punto del calendario donde la guía cambia de versión.

\subsection{Política de actualización}

\begin{uxtabla}{|L{3.6cm}|Y|}{Política de actualización de la guía}
  \uxheadrow \thd{Regla} & \thd{Contenido} \\
  Vía única de cambio & Se edita la hoja de estilo o el generador, se vuelve a emitir y se compila el documento. Ninguna de las cuatro salidas se edita a mano \\
  Verificación obligatoria & Idempotencia de las cuatro salidas, invariante de las once anclas, ausencia de color escrito a mano y coincidencia de la versión en los tres sitios \\
  Compatibilidad & Ningún nombre de token se renombra ni cambia de valor dentro de una versión mayor; los nombres heredados se conservan como alias derivados de la escala \\
  Numeración & La cifra menor sube cuando se añaden tokens o reglas sin romper nombres; la mayor, cuando un nombre cambia de significado \\
  Cadencia & La guía se versiona con el producto, en la revisión de la entrega semanal \\
\end{uxtabla}

\subsection{Entradas abiertas de la versión 1.1}

Dos, con su condición de entrada escrita: \textbf{densidad compacta}, que entra cuando existe una segunda escala de altura de fila verificada sobre la tabla de datos; y las \textbf{correcciones que salgan de la prueba de usabilidad de la Actividad 5}, que entran con el número de participantes y la medición que las motiva. Ninguna de las dos se anticipa aquí con valores: una entrada de bitácora sin dato verificado es exactamente el tipo de afirmación que esta guía existe para impedir.
